\scanpage{34}
\mspage{16}
Cela signifie qu'on se donne, pour $\alpha$ fixé,

1°) un graphe \inserted{connexe} $(S_{(\alpha)},\widetilde A_{(\alpha)},\deleted{\unsure{\widetilde A_{(\alpha)}}}\sigma_{\widetilde A_{(\alpha)}},\widetilde A_{(\alpha)}\to S_{(\alpha)})$

2°) on \deleted{\unsure{se}} donne une application ``genre''
\[
S_{(\alpha)}\xrightarrow{\unsure{\gamma_\alpha|g_\alpha}}\mathbb N
\]

3°) on donne une application $I_\alpha\to S_{(\alpha)}$, i.e. on ``répartit'' les \deleted{\unsure{divers}} ``pts marqués'' au dessus de $\alpha$ \inserted{(i.e. les éléments de $I_\alpha$)} entre les sommets $\beta\in S_{(\alpha)}$ de $G_{(\alpha)}$

4°) on donne une application $\widetilde A_\alpha\to S_{(\alpha)}$, i.e. pour $\forall$ \deleted{flèche}\inserted{arc} du graphe $G$, d'origine $\alpha$, on \deleted{\unsure{se donne}} choisit une \inserted{nouvelle} ``origine'' parmi les sommets $\beta\in S_{(\alpha)}$ du graphe $G_{(\alpha)}$

5°) Les données 1°) : 4°) définissent une maquette, \inserted{0-connexe, $G_{(\alpha)}$} et on prend soin que celle-ci soit un MD-graphe, \deleted{\unsure{ce}} ce qui signifie ici pour pour $\forall\beta\in S_{(\alpha)}$ ``de genre 0'' i.e. tel que $\unsure{\gamma_\alpha|g_\alpha}(\beta)=0$, on ait $\operatorname{card}I_{\alpha,\beta}+\operatorname{card}\widetilde A_{\alpha,\beta}\geq 3$.

Ceci étant, on trouve un nouveau \inserted{MD-}graphe\inserted{$G'$}, \inserted{qui sera} spécialisé du \inserted{MD-}graphe $G$.

Je voudrais préciser de façon symétrique la relation entre les MD-graphes $G$ et $G'$, et j'introduis à cet effet une digression sur
\mspage{33}
